\section{Application Details and Additional Comparisons}
\label{app:applications}

This appendix supports Section~\ref{sec:applications}.
It preserves the full PCS and Flock methodology, additional comparisons,
and the distinction between measurements and projected costs.

\subsection{Silent OT Integration}
\label{app:pcg-details}
\paragraph{Regular-Noise Batches.}
The regular-noise integration reuses the full precomputed rate-$1/2$ BCH--IMT code,
with $(t,s)=(128,19)$ and weight-five feedback. There is no heuristic
permutation bank or per-batch code refresh. Each batch expands fresh trees,
compresses their outputs, and hashes both sender messages.
Table~\ref{tab:spin-regular-current} reports sender computation
at two batch sizes. All entries use VAES and streaming output stores.

\begin{table}[!htbp]
\centering
\caption{Regular-noise sender computation. Initial setup,
base-correlation generation, and network transport are excluded.}
\label{tab:spin-regular-current}
\begin{tabular}{@{}lrr@{}}
\toprule
Timed work & $K=2^{18}$ & $K=2^{20}$ \\
\midrule
Tree expansion (ms) & 1.033 & 4.093 \\
SPIN compression (ms) & 1.482 & 11.351 \\
Hash both sender messages (ms) & 0.428 & 1.728 \\
Total (ms) & 2.944 & 17.186 \\
Million hashed OTs/s & 89.0 & 61.0 \\
\bottomrule
\end{tabular}
\end{table}

Both sizes use $128$ noise partitions, with $4096$ and $16384$ leaves
per tree, respectively, and no padding. Each batch consumes $1536$ or
$1792$ fresh base-OT correlations, supplied outside timing.
The noise configurator uses an effective-distance tuning value of $0.25$;
this is not a certified relative distance or a new hardness claim.
The runs use the hardware, compilers, and three-process measurement protocol
described below for stationary mode. Setup, workspace and output allocation,
base-correlation generation, transport, and output consumption are untimed.
The measured calls include routine allocation and local message enqueueing.
Separate paired runs check OT correctness.

With ordinary cached stores, totals are $3.066$ and $17.495$\,ms.
Including an immediate sequential scan of the streaming outputs raises
the totals to $3.171$ and $17.971$\,ms. A stationary run matched to these
measurements takes $3.103$\,ms at $K=2^{18}$, or $84.5$ million hashed OTs
per second; Table~\ref{tab:spin-stationary-ot}, measured separately,
reports $3.101$\,ms.
The modes consume different base correlations, whose generation is excluded;
these timings do not establish which complete protocol has lower cost.

\paragraph{Stationary Batches.}
The stationary integration uses the rate-$1/2$ BCH--IMT encoder at $K=2^{18}$,
with $(t,s)=(128,19)$ and weight-five feedback.
The parties retain a precomputed permutation bank and tree seeds between
batches. Each batch refreshes the routing parameters and IMT masks, generates
fresh leaves, and applies the transposed encoder. The refreshed permutations
form a heuristic family; the uniform-setup distance certificates do not apply.

\begin{table}[!htbp]
\centering
\caption{Stationary sender computation for $2^{18}$ hashed OTs per batch.
Both modes use VAES; only the output-store policy differs.
Initial setup, base-correlation generation, and network transport are excluded.}
\label{tab:spin-stationary-ot}
\begin{tabular}{@{}lrr@{}}
\toprule
Timed work & Cached stores (ms) & Streaming stores (ms) \\
\midrule
Fresh leaves and padding & 0.941 & 0.917 \\
Code refresh and compression & 2.213 & 1.752 \\
Hash both sender messages & 0.423 & 0.429 \\
Total & 3.577 & 3.101 \\
\bottomrule
\end{tabular}
\end{table}

The stationary configuration uses $312$ trees of $1680$ leaves and pads the
remaining $128$ coordinates to length $2^{19}$ with zeros. Cached leaf seeds
occupy about $8$\,MiB per party. Each batch consumes $312$ fresh base-VOLE
correlations, supplied outside timing. Initial tree expansion, permutation-bank
construction, and workspace and output allocation are also untimed.
The timer covers fresh-leaf generation and sums, padding, code refresh,
compression, and AES hashing of both sender messages.
Routine allocation and local message enqueueing within these calls remain
timed; transport, transcript draining, and output consumption do not.

For both modes in Table~\ref{tab:spin-stationary-ot}, we report the median
of three process medians.
Each process measures $101$ batches after five warmups. Stage medians need not sum to
the median total. The processes run serially on one Ryzen 9 7950X core
at a requested $4.5$\,GHz with boost disabled. libOTe and cryptoTools use
GCC~13.4 with VAES enabled; the SPIN library uses GCC~15.2.0 with Zen~4 tuning
and runtime-dispatched AVX-512 BCH kernels.
Before timing, the harness initializes both endpoints and checks their OT
outputs. It then releases the receiver's buffers to measure the sender alone.
Separate paired runs check agreement across successive code refreshes.

Streaming stores avoid filling the cache with the $8$\,MiB sender-output
buffer. Their main benefit appears in the next batch's compression, not
in hashing itself. Including an immediate sequential scan of all output bytes
raises the streaming total to $3.425$\,ms, or $76.5$ million OTs per second;
cached stores take $3.884$\,ms with the same scan.
This scan is a memory-access control, not a downstream protocol benchmark.
Streaming stores remain opt-in because the benefit depends on batch size
and the consumer's access pattern.

\paragraph{Regular-Noise Integration at $K=2^{20}$.}
We also report an integration measurement with a different length, noise
configuration, and encoder implementation. It is not a matched baseline for
Table~\ref{tab:spin-stationary-ot}, and its sender totals are not comparable
with Table~\ref{tab:spin-regular-current}.
We replace the linear map in libOTe's semi-honest Silent OT path with the
rate-$1/2$ BCH--IMT encoder at $K=2^{20}$, $(t,s)=(128,19)$ and weight-five
feedback. Table~\ref{tab:spin-ot} reports protocol computation, not an
encoder-based projection.

\begin{table}[!htbp]
\centering
\caption{Measured time per party for $2^{20}$ outputs in the Silent OT
integration. Base OT, setup, buffer allocation and network latency are excluded.}
\label{tab:spin-ot}
\begin{tabular}{@{}lrr@{}}
\toprule
Timed work & Sender (ms) & Receiver (ms) \\
\midrule
Expansion + SPIN & 19.531 & 20.625 \\
Expansion + SPIN + hashing & 22.099 & 21.841 \\
\bottomrule
\end{tabular}
\end{table}

The regular binary noise has $400$ partitions of $5242$ coordinates each;
the remaining $352$ coordinates of the $2^{21}$-coordinate input are zero.
The configuration subtracts those unused coordinates from the distance
input to libOTe's noise-weight calculation. This benchmark specifies the
implementation and its costs; it does not supply a new hardness claim for
the resulting structured LPN instance.

The harness installs synthetic matching base correlations, constructs the
public SPIN setup, and allocates and touches both endpoints' buffers before
timing. It runs the sender and receiver serially on the same core, using
the sender's actual buffered transcript for the receiver. Each endpoint
has its own timer and workspace. The measured calls perform expansion
and the in-place linear map; the receiver also packs its choice bits.
Transcript handoff, network latency, output checks and object destruction
are excluded. Online communication is $179{,}624$ bytes, excluding base
correlations.

Without hashing, the receiver stores its choices in the low bits of its
blocks; the correlated-OT relation holds in the remaining $127$ bits.
With hashing, the harness checks each selected OT message against the
corresponding sender output. All $2^{20}$ outputs pass in every trial,
including warmups, in both modes. Independent dense-oracle encoder tests
also pass at $K=2^{16},2^{18},2^{20}$.

The runs use one Ryzen 9 7950X core, pinned to CPU~15 at a requested
$4.5$\,GHz with boost disabled, and GCC~15.2.0 with
\texttt{-O3 -march=znver4}. Each entry is the median of nine calls after
one warmup. Hash-enabled totals are medians of per-call sums of protocol
and hashing time, not sums of independently computed medians.
The no-hash sender times range from $19.474$ to $19.911$\,ms and receiver
times from $20.532$ to $20.754$\,ms.

For context, the chosen-block BAA evaluation reports $22.810$\,ms for Golay
and $24.278$\,ms for RM for the same $K$ and rate, using an Intel
i9-14900HX and MSVC. Table~\ref{tab:transposed-comparison} instead uses
our same-host measurements and does not combine cross-machine timings.


\subsection{Polynomial Commitment Details}
\label{app:pcs-details}

We use SPIN as the row encoder in a Brakedown-style polynomial commitment
scheme~\cite{brakedown23}. This application tests whether fast ordinary
encoding translates into fast commitment and opening. The implementation
has two principal changes: it instantiates the code with SPIN, and it
organizes the computation around packed binary data. At $512$\,MiB of
input, the resulting prover commits and opens in $509$\,ms on one core.

\paragraph{Construction.}
Let a Boolean evaluation table define a multilinear polynomial $f$ over
$\mathbb F:=\F_{2^{128}}$. A public index mapping arranges its $RK$ bits
into a matrix $M\in\F_2^{R\times K}$. The prover applies the same SPIN
encoder $\operatorname{Enc}$ to each row $M_i$, obtaining
$U\in\F_2^{R\times N}$ with $N=2K$, and Merkle-commits to its columns.
We also write $\operatorname{Enc}$ for the extension of its binary generator
to $\mathbb F$.
For a requested evaluation point $z$, let $\beta\in\mathbb F^R$ and
$\gamma\in\mathbb F^K$ be the Boolean Lagrange weights for the row and
column variables under this mapping. Thus
$f(z)=\sum_{i,j}\beta_i\gamma_jM_{ij}$.

In the interactive protocol, the commitment, $z$, and claimed value $v$
are fixed before the verifier samples $\alpha\gets\mathbb F^R$ and sends
it to the prover. The honest prover returns the explicit folded messages
\[
  m_T:=\sum_{i=1}^R\alpha_i M_i,
  \qquad
  m_E:=\sum_{i=1}^R\beta_i M_i
  \qquad\text{in }\mathbb F^K.
\]
After receiving both messages, the verifier draws $q$ independent uniform
column indices from $[N]$ with replacement and requests openings at the distinct
sampled indices. The prover returns those columns and their Merkle
authentication. The verifier rejects malformed messages or invalid
authentication. It also checks $\sum_j\gamma_j(m_E)_j=v$ and, for every
sampled column $j$,
\[
  \operatorname{Enc}(m_T)_j=\sum_i\alpha_iU_{ij},
  \qquad
  \operatorname{Enc}(m_E)_j=\sum_i\beta_iU_{ij}.
\]
It accepts only if all checks pass. Both folds share the query set;
deduplication does not change the event of missing all inconsistent columns.

The measured implementation uses domain-separated BLAKE3 transcript
challenges in place of verifier randomness. It absorbs the parameter
identifier, commitment root, point, and claimed value before deriving
$\alpha$. It then absorbs $m_T$ followed by $m_E$ before deriving the
column indices. The verifier reproduces this order.

\paragraph{Making the Binary Code Useful.}
Expanding every bit to an extension-field element would erase much of
SPIN's encoding advantage. We therefore keep the matrix packed and
encode bounded groups of rows. The internal layout follows the input
index mapping, avoiding a full-matrix bit transpose. Batched column
processing and table-based row combinations reduce the data movement
and arithmetic around the encoder. Workspaces are allocated before
repeated proving. These changes make encoding and column hashing account
for approximately $79\%$ of the fastest standalone prover time.

\paragraph{Parameters.}
We condition once on the selected code satisfying the distance bound
of Theorem~\ref{thm:finite-bch-spin}. Extending its binary generator to
$\mathbb F$ preserves minimum distance. At $K=2^{16}$ and $2^{18}$,
we use $2045$ column samples and one random testing fold. For a complete
matrix fixed before the testing challenge, the generic-code proximity
calculation~\cite{diamondPosen24} gives an error below $2^{-100}$ for the
interactive matrix test and evaluation check. We call this
the conditional 100-bit fixed-matrix testing target. It excludes the one-time setup
event and the separate commitment-extraction and Fiat--Shamir reductions;
the calculation does not establish a complete composed security theorem.
Our verification code checks this bound and its parameters with exact arithmetic.

\paragraph{Standalone Performance.}
Table~\ref{tab:spin-pcs-standalone} compares SPIN--Brakedown with
Ligerito~\cite{ligerito25} on the same $2^{32}$-bit input. SPIN uses
$R=K=2^{16}$ for the square shape. The longer-row shape has $R=2^{14}$
and $K=2^{18}$. Both encode to $1$\,GiB.
Longer rows reduce the SPIN proof
size by $31.6\%$ for a $3.1\%$ increase in prover time. The square
configuration spends $221$\,ms encoding, $180$\,ms hashing columns,
and approximately $96$\,ms computing the two row combinations.

All measurements use one worker on the Ryzen 9 7950X, pinned to CPU~15,
with a requested $4.5$\,GHz frequency and boost disabled. Each SPIN entry is
the pooled median of ten trials from two processes, each excluding one
warmup. For SPIN, setup, workspace allocation, outer serialization, and independent
evaluation for validation are outside the prover timer. Verification
excludes decoding and workspace allocation. Every timed proof verifies.

Ligerito's rows use the optimized implementation that also supplies the
Flock baseline in Section~\ref{sec:spin-flock}. Its ring-switching layer
opens the binary multilinear polynomial at an ordinary evaluation point;
no Flock computation or cached partial evaluation is supplied. Fast100
and Slim100 use initial code rates $1/2$ and $1/4$, respectively, and
their shipped query and grinding parameters. Their timings pool ten trials
from two processes per profile, each with one excluded warmup. Input
preparation, statement serialization, independent evaluation, outer proof
serialization, and decoding are excluded; internal allocations remain
timed. These profiles give much smaller proofs and faster verification,
at higher prover cost.

\paragraph{Comparison with Bolt.}
\label{app:bolt-estimate}
At the same $512$\,MiB input volume and on one core, the fastest measured
Bolt-max commitment takes $1342$\,ms, compared with $401$\,ms for
SPIN--Brakedown: a $3.34\times$ difference. Both are measured on the
same host. Bolt~\cite{bolt26} encodes
$2^{27}$ elements of $\F_{2^{32}}$; our input has $2^{32}$ Boolean entries.
The comparison therefore matches input bytes, rather than polynomial
dimension or coefficient field. It includes each implementation's code,
layout, and commitment processing. The authors' public Bolt implementation
provides commitment but does not implement the complete opening prover.
We optimized Bolt's x86 expander by specializing its $128$-symbol columns
and replacing each heap accumulator with a fixed-size array.
The compiler keeps the accumulator in SIMD registers; no handwritten
assembly is used. Against the unmodified baseline, expander time falls
from $733$ to $331$\,ms and commitment time from $1757$ to $1342$\,ms.
The graph, parameters, Reed--Solomon kernel, and hashing are unchanged;
baseline and optimized commitment roots match. Timings are medians of
six measured trials per binary, across two processes with one excluded
warmup each. The repository records the source hashes and timing
boundaries\repositorycite.

\paragraph{Opening Cost in Isolation.}
Table~\ref{tab:pcs-opening} separates opening from commitment. The two
SPIN configurations both open in approximately $108$\,ms, including the
row combinations and authenticated column queries. These are measured
opening times from the same trials as Table~\ref{tab:spin-pcs-standalone};
they exclude outer serialization. Sending the folded messages explicitly
keeps the prover's opening work small, at the proof sizes reported above.

For Bolt, we measured the dominant opening computations. We did not
implement the full opening; Table~\ref{tab:pcs-opening} reports a calibrated
projection.

\begin{table}[tb]
\centering
\caption{One opening after commitment, at $512$\,MiB input volume on one
core. Bolt entries are calibrated projections, not complete prover
measurements or runtime bounds.}
\label{tab:pcs-opening}
\small
\begin{tabular}{@{}llr@{}}
\toprule
Configuration & Evidence & Opening (ms) \\
\midrule
% Generated by build_application_tables.py; data SHA-256 07c6005a5a76085c42de1fb37977c1d8b2633a7508deef3ce9d56b95db534297
% Bolt projection SHA-256 f7506e7c23e0b4831398fae38c2006de346a212140b034e22ded2ab51c9732cb
% Ligerito data SHA-256 ddf993619da6426cae615e196bc8485742539c188b7def0bf368feb94a181839
SPIN--Brakedown (square) & Measured & 108 \\
SPIN--Brakedown (longer rows) & Measured & 108 \\
Ligerito (Fast100) & Measured & 3389 \\
Ligerito (Slim100) & Measured & 3668 \\
Bolt-max (amortized limit) & Calibrated projection & 1273 \\
Bolt-max (non-amortized) & Calibrated projection & 1729 \\
\bottomrule

\end{tabular}
\end{table}

Bolt's opening additionally uses inner Reed--Solomon proofs and a proof
of the sparse-matrix relation~\cite{bolt26}. To estimate its cost, we
measure row evaluation, a random column fold, and sumchecks, and use
complete Ligerito proximity runs as proxies for the inner proofs. The
calibrated projection is $1273$\,ms in the amortized limit, which assigns
negligible cost per opening to the matrix-relation proof. Adding a
calibration of that proof's leading multiplication work gives $1729$\,ms
without amortization. The inner proxies differ from Bolt's constrained
and batched proofs; basis preparation, some matrix-proof work, outer
authentication, and transcript integration remain unmodeled. The
repository records these choices and the measured components\repositorycite.
These figures are the projected cost of the selected implementation
approach, not a lower bound on Bolt's opening time.

\paragraph{Estimated Bolt Communication.}
The Bolt-max row in Table~\ref{tab:spin-pcs-standalone} uses the same
amortized-limit model for time and communication. Adding the measured
commitment to the projected opening gives $2614$\,ms, using unrounded inputs.
For communication, we use the paper's nominal 100-bit query formulas,
with parity-check distance $\gamma=0.013$ and rate-$1/2$ Reed--Solomon
encoding. They give $15{,}962$ message queries and $241$ syndrome queries
\cite[Section 7.3]{bolt26}.

The measured commitment has $2^{20}$ message columns and $2^{18}$
syndrome columns, each containing $128$ four-byte symbols.
Uniform distinct queries within each part transmit $7.912$\,MiB of
column data and an expected $2.600$\,MiB of batched Merkle authentication.
The authentication calculation includes the two dummy sibling hashes in
the implementation's non-power-of-two tree.
We add $451{,}184$ and $177{,}200$ bytes from the two verified inner-proof
proxies used for timing, $2048$ bytes of row evaluations, $1920$ bytes for
two sumcheck transcripts, and $128$ bytes of scalar-message allowance.
The resulting communication estimate is $11.115$\,MiB, rounded to $11.1$.

The inner proxies are complete Ligerito proximity proofs, not Bolt's
constrained proofs. They include their own commitments and do not exploit
Bolt's virtual syndrome oracle or batching. The estimate assigns zero
per-opening communication to the shared matrix proof in the amortized
limit; finite batches must pay their share. It excludes container framing
and the original commitment root. These substitutions make the total a
communication model, not a generated Bolt proof.
As a secondary comparison, keeping the unresolved overhead in Bolt's
published $4$\,GiB proof gives $10.9$--$11.5$\,MiB at $512$\,MiB,
depending on whether Bolt's MB denotes $10^6$ or $2^{20}$ bytes. That
agreement does not independently validate the inner-proof model. The
repository includes the calculation and its tests\repositorycite.

\paragraph{Compatibility with Blaze.}
SPIN is also compatible with Blaze's generic code-switching
framework~\cite{blaze25}. Given the stated distance bound, its linear
encoder can be combined with a generic interactive oracle proof of
proximity for the encoding and multilinear-evaluation relation. This
gives a theoretical SPIN--Blaze composition with asymptotically smaller
proofs than our Brakedown-style implementation. The inner proof replaces
the explicit folded messages, allowing longer rows and shorter
authenticated columns. The large proofs in our measurements therefore
come from the chosen PCS construction. We give no concrete SPIN--Blaze
instantiation, parameters, or performance.

The opening comparison motivates testing the complete application:
commitment speed alone does not determine the cost of a PCS inside a
proof system. Section~\ref{sec:spin-flock} measures both commitment and
opening after integration into Flock, where the claim interface and
available intermediate computations also affect their cost.

\subsection{Flock Integration and Full Results}
\label{app:flock-details}

We integrate SPIN--Brakedown into Flock~\cite{flock26} to measure its
effect inside a larger prover. Flock checks batches of Boolean constraints
and reduces them to claims about a committed witness. Our integration
replaces its Ligerito PCS~\cite{ligerito25} with the construction in
Section~\ref{sec:spin-pcs}. SPIN improves prover throughput
by $2.29\times$ and $1.63\times$ at the two measured workloads.

\paragraph{Adapting the Opening Interface.}
Flock's final claims carry interpolation weights on a small group
of witness coordinates. We incorporate these weights into the PCS row
combination, so the opening proves the functional that Flock actually
requests. A public permutation of witness variables preserves the six-coordinate
interpolation prefix and groups known padding into zero rows. The commitment
omits these rows; the opening accounts for their zero contribution. Flock
generates the packed witness directly in this compact layout, avoiding a
separate permutation pass. The adapter uses
two separate openings with $2110$ column samples each and $K=2^{16}$.
Their conditional fixed-matrix error bounds sum to less than $2^{-102}$.
The benchmark proves the committed batch's
compression constraints; binding public hash inputs and outputs is a
separate application wrapper.

\paragraph{Improving the Surrounding Prover.}
During integration, we found that field arithmetic, repeated witness
passes, and temporary storage could obscure the PCS improvement. We
optimized Flock's x86 kernels, reused workspaces, and evaluated the
fixed linear constraints directly. The Ligerito configuration receives
the applicable improvements as well. In particular, its existing
column buffers can supply an intermediate zerocheck value, avoiding
duplicate work. These changes preserve the checked protocol messages.
The comparison below uses both PCS backends within this optimized
implementation, rather than attributing shared improvements to SPIN.

\begin{table}[tb]
\centering
\caption{Flock proofs of BLAKE3-compression constraints on one Ryzen
core. SPIN and Ligerito, using its Fast100 profile, are measured with the
applicable shared optimizations. Bolt entries are optimistic projections assuming free claim
conversion, shared opening work, and amortized matrix-proof cost;
Bolt--Flock is not implemented. Dashes denote unestimated quantities.}
\label{tab:spin-flock}
\small
\begin{tabular}{@{}rlrrr@{}}
\toprule
Compressions & PCS & Prover (ms) & Verify (ms) & Proof (MiB) \\
\midrule
% Generated by build_application_tables.py; data SHA-256 07c6005a5a76085c42de1fb37977c1d8b2633a7508deef3ce9d56b95db534297
% Bolt projection SHA-256 f7506e7c23e0b4831398fae38c2006de346a212140b034e22ded2ab51c9732cb
% Bolt x86 commitment SHA-256 caf348da05c88c1ce12b00dd955a7e6db46b892512b9732dafabf7f64360ed22
$16,384$ & SPIN--Brakedown & 94 & 37 & 6.2 \\
$16,384$ & Ligerito & 215 & 28 & 0.3 \\
$16,384$ & Bolt (projection) & 228 & --- & --- \\
$65,536$ & SPIN--Brakedown & 351 & 42 & 10.5 \\
$65,536$ & Ligerito & 571 & 28 & 0.4 \\
$65,536$ & Bolt (projection) & 967 & --- & --- \\
\bottomrule

\end{tabular}
\end{table}

\paragraph{Measurement Method.}
The SPIN and Ligerito configurations run in one executable with the same deterministic
input stream and compiler settings. They use the single-worker hardware
environment of the standalone PCS experiment: one Ryzen 9 7950X core at
$4.5$\,GHz, with boost disabled. The build uses Rust 1.96.0, GCC 15.2,
native CPU instructions, and 64-byte function alignment. Each workload has
two processes per backend in a balanced serial order. Each process excludes
five warmup and settling proofs and measures four trials; each entry is
the mean of two process medians. Every generated proof is verified.
Prover time includes witness generation, commitment, and proving. Setup,
verification, and outer proof serialization are outside the interval.
Serialization inside the SPIN opening remains timed because its bytes
enter the protocol transcript.

\paragraph{Prover Costs.}
At the smaller workload, commitment takes $15.7$\,ms with SPIN--Brakedown
and $43.7$\,ms with Ligerito; opening takes $17.1$ and $102.8$\,ms.
At the larger workload, these costs are $62.2$ versus $162.0$\,ms for
commitment and $46.0$ versus $141.7$\,ms for opening. The larger-workload
PCS savings account for about $195$\,ms of the $221$\,ms total reduction.
Encoding alone takes $7.9$ and $28.0$\,ms with SPIN, versus $26.9$ and
$98.6$\,ms with Reed--Solomon. These are complete integrations with
different layouts and hashing; the measurements do not isolate an
encoder-only substitution.

Four times the workload raises SPIN commitment and opening from $32.7$
to $108.2$\,ms, improving their cost per compression. The surrounding
prover takes $242.7$\,ms at the larger size, about $69\%$ of the total.
Ligerito opening grows only from $102.8$ to $141.7$\,ms, which helps
explain the smaller relative prover speedup at that size. Within SPIN
commitment, leaf hashing grows from $6.7$ to $33.1$\,ms and is a
remaining optimization target.

\paragraph{An Optimistic Bolt Projection.}
Table~\ref{tab:spin-flock} also considers replacing the PCS with
Bolt, using the calibrated opening costs from Appendix~\ref{app:pcs-details}.
We give Bolt favorable assumptions: free conversion of Flock's weighted
claims, shared opening work across its two claims, and negligible
matrix-proof cost per opening in the amortized limit. We keep the
surrounding prover time measured with SPIN and substitute Bolt's measured
commitment and projected opening costs. The resulting totals are
$228$ and $967$\,ms, compared with the measured SPIN totals of $94$
and $351$\,ms. Thus the projected Bolt configuration remains slower even
with these allowances. We did not implement Bolt--Flock; the projection
describes this substitution and is not a runtime bound.

The SPIN--Brakedown integration favors prover time at a substantial
communication cost. At the larger workload, its proof is $10.5$\,MiB,
compared with $0.4$\,MiB for Ligerito, which also verifies faster.

\FloatBarrier
